\uuid{bJIa}
\titre{Raisonnements}
\auteur{bodin ; croizat ; sacre}
\organisation{exo7}
\licence{CC BY-NC-SA 4.0}
\datecreate{2022-03-08}
\niveau{L1}
\module{Algèbre}
\chapitre{Logique, ensemble, raisonnement}
\sousChapitre{Absurde et contraposée}
\difficulte{2}
\theme{Raisonnements}
\exo7qcmid{hilisit.logique.048}

\contenu{
\begin{question}
\qclassification{100.03}
Pour montrer que les solutions réelles de l'équation $\vert x+1\vert =2$ sont $1$ et $-3$, on peut utiliser :
\begin{answers}
	\bad{Un raisonnement par contraposition.}
	\good{Un raisonnement par disjonction des cas.}
	\bad{Un raisonnement par l'absurde.}
	\bad{Un raisonnement par récurrence.}
\end{answers}
\begin{explanations}
Soit $x\in\Rr$. On distingue les cas $x+1\ge 0$ et $x+1<0$ :

Si $x+1\ge 0$, alors $\vert x+1\vert =x+1$ et l'équation devient $x+1=2$, qui a $1$ pour unique solution.

Si $x+1<0$, alors $x+1=-(x+1)$ et l'équation devient $-(x+1)=2$, qui a $-3$ pour unique solution.
\end{explanations}
\end{question}
}
